\section{文化与思想：天下失去中心后，人们怎样重新解释秩序}

战国的思想不在战场之外。城池反复攻守，诸侯以赋役和军功求强，旧有的王室名分又不足以裁断新王的竞争；在这种处境中，礼、利、法、兵、天与人的关系，都成了可以向君主进言、也可以彼此反驳的问题。所谓“诸子”，不是一套整齐的同代名册：有些是跨国活动的人，有些是师承网络，有些是后来由篇章、传记和目录归拢出来的传统。

知识也并未只在大都城发生。前433年曾侯乙墓中的编钟、礼器及其铭文，留下了可与传世叙事互校的近同时物证；\eventanchor{WQ-0433-ZENGHOUYI}{战国·曾侯乙墓}让我们看见礼乐、冶铸和楚系诸侯关系的实际物质条件，却不能由一座墓推演整个战国社会的音乐制度。约前323年中山王墓的列鼎、车马器和刻铭铜器同样显示，小国称王时也须用可见的礼器组织王室与贵族的认同，见\eventanchor{WQ-0323-ZHONGSHAN-ROYAL-CEMETERY}{战国·中山王墓礼器}。礼不是已经失效的古老外壳；它仍是各国在资源、工艺和名分竞争中使用的语言。

\subsection{人随道路与朝廷流动，主张也随之改变位置}

早期的西河已经说明，讲学并不等于远离政事。约前403至前390年，子夏、田子方等士人在魏的西河讲学，魏廷以宾礼吸纳儒者和军政人才；\eventanchor{WQ-0403-ZIXIA-XIHE}{战国·西河讲学}所能确知的是一种把文教、选贤与边地国家建设连在一起的网络，而不是一份固定课程表。魏、齐、楚、秦之间的道路与使聘，使士人得以寻找新的依附者；国君则以馆舍、俸禄、相位或可被倾听的机会争取这些人。

这是一种双向选择，而非学者单方面教化君主。孟子游梁的年代和对话的编次仍有争议，\eventanchor{WQ-0350-MENCIUS-WEI}{战国·孟子游梁}与\eventanchor{WQ-0340-MENCIUS-WANGDAO}{战国·孟子论王道}所保存的，是魏廷在富国强兵的焦虑中听到“利义”“养民”这一套反驳语言；《孟子》的问答体不能改写成梁惠王朝议的逐字记录。约前331年“何必曰利”的著名开场，见\eventanchor{WQ-0331-MENCIUS-LIANGHUI}{战国·孟子见梁惠王}，其力量正在于把粮食、赋税和战争的计算重新放进君主应如何对待百姓的问题中，而不在于它已经成为魏国的国策。

齐廷提供的是另一种通道。孟子在齐的活动及其同齐宣王的讨论，可由\eventanchor{WQ-0330-MENCIUS-QI}{战国·孟子至齐}和\eventanchor{WQ-0321-MENCIUS-QI-COUNSEL}{战国·孟子游齐进言}追踪：王道、民生与用人因而进入扩张中的强国宫廷，但君主是否采纳、采纳到何种程度，不能由篇章中的一席话代定。羊易牛的故事，\eventanchor{WQ-0320-QI-XUAN-OX-EXCHANGE}{战国·齐宣王易牛}，把祭祀牲畜、恻隐之心和礼制成本并置；它是后世不断解释的思想材料，不是一则可以核验的宫廷账簿。儒者的“仁”“礼”在此更像衡量征发、用兵和君臣责任的竞争性尺度，而不是任何国家稳定、一致的行政方案。

齐的稷下尤其不宜被想象成一所制度完备的近代学校。约前四世纪中叶以后，齐廷向游士提供居处、名号和议论空间的惯例发展，见\eventref{WQ-0347-QI-JIXIA-PATRONAGE}{战国·稷下延士}；约前320年齐宣王延纳学者，见\eventref{WQ-0320-QI-XUAN-JIXIA}{战国·宣王延士}。传世材料足以说明临淄把财富转化为文化声望和人才市场，不能让我们列出校舍、经费、学籍或所有成员的同时在场。荀子约前285至前260年游学稷下并被后世称为祭酒，\eventanchor{WQ-0285-XUNZI-JIXIA}{战国·荀子游学稷下}；其礼、法与国家秩序的系统论述，显示同一知识场域可以容纳彼此不相同的答案，并非“稷下学派”的统一教义。邹衍等人在稷下的活动，\eventanchor{WQ-0300-QI-JIXIA-YANZHOU}{战国·邹衍与稷下}，又使五德终始一类王朝更替语言可以服务于列国的合法性想象。齐国失去临淄时，这个网络亦随都城的税入、保护和声望受挫，见\eventref{WQ-0284-YAN-CAPTURE-LINZI}{战国·燕军入临淄}；文化资本不能替代城防、盟友和粮道。

齐人王蠋在燕军占领时拒绝合作而死的故事，\eventanchor{WQ-0283-QI-WANGZHU-DIES}{战国·王蠋拒燕而死}，后来被《史记》《战国策》写成士人忠义的鲜明见证。它可帮助理解占领者为何需要争取齐地士人与城邑精英，也说明田单复国叙事如何获得道德正当性；但一位忠义者的传记不能替代齐地民众、地方势力在战乱中的复杂选择。

燕昭王“黄金台”招贤的故事同样须置于记忆与政策之间。它以\eventanchor{WQ-0311-YAN-GOLDEN-TERRACE}{战国·黄金台招贤}概括燕廷对外来士人的礼遇承诺，能说明求贤形象为何长期成为燕国由弱转强的叙述资源；台址、赏赐和来者名单却不能据《战国策》的修辞逐项坐实。惠施在魏的名辩与合纵主张，也显示辩说可以进入朝廷的战略困境，见\eventref{WQ-0319-HUISHI-WEI}{战国·惠施在魏}；它并不证明一篇辩论就足以替魏弥补马陵之后的兵力与地理劣势。

\subsection{几种传统回答的，不是同一道考题}

后世称为儒、墨、道、法的传统，确有彼此可辨的关怀，却不应被压成四个同时覆盖全国的政党。儒者以仁政、礼与养民批评只问功利的国家竞争；墨者则把兼爱、非攻、节用与可操作的守城知识放在一起。约前四世纪前期，墨者以巨子、弟子网络传授器械、守城和游说技能，见\eventanchor{WQ-0390-MOZI-ORGANIZATION}{战国·墨者组织与传习}。这类材料说明战争使可复制的工程知识成为跨国资源，也说明技术训练与反战论述可以交织；《墨子》各篇的成书层次复杂，不能逐条充作某一年某城的守备实录。

道家文本提供的并非对政治毫无兴趣的退场姿态。约前四世纪末，庄周及相关篇章在宋、魏、齐之间流传，质疑功名依附与强制秩序，见\eventanchor{WQ-0300-ZHUANGZI-CIRCULATION}{战国·庄周与道家文本流传}。从《庄子》与后出的传记可以把握一种对战国权力、生计与名分的反思，却不能把每一则寓言还原为庄周在某国朝廷的亲历。它和儒、墨的关系也不必预设为整齐对垒：文本、弟子和读者跨越的道路，比后世的门派边界更不规整。

所谓法家尤其不是秦国的专利。魏、韩、楚、秦都曾在强邻竞争中讨论考核、赏罚、军功与官吏控制；例如\eventref{WQ-0354-SHENBUHAI}{战国·申不害相韩}与\eventref{WQ-0356-SHANGYANG}{战国·商鞅变法}，分别留下“术”与“法”如何被后人整理的线索。韩非约前260至前234年撰述《孤愤》《五蠹》等篇，\eventref{WQ-0260-HANFEI-WRITING}{战国·韩非撰述法术思想}；将法、术、势编织得更为系统，是在韩国受秦挤压的处境中形成的理论回应，不可倒推为韩国已实行的一套完整政策。

前234至前233年，秦廷读韩非著作并召其入秦，韩国则试图以公子身份和思想声望取得缓冲，见\eventref{WQ-0234-HANFEI-QIN-RECEPTION}{战国·秦廷召韩非}。韩非旋即死于秦，\eventanchor{WQ-0233-HANFEI}{战国·韩非入秦而死}；《史记》将李斯、猜忌与构陷收束为强烈的传记转折，细部仍有塑形。较稳妥的结论是，能被强国阅读的文本并不能替弱国改变吞并战略；理论被吸收、改造或后来重读，也不等于它已经原样成为秦廷的统一政策。

\begin{textwindow}[从“说服”读到什么]
《孟子》《庄子》《韩非子》与《战国策》保存了战国人怎样把政治难题说成一段可传诵的言辞。它们让后人听见王道、功利、名辩与法术的冲突，但不是录音式的朝议纪录。阅读这些文本时，应同时追问：谁在何种成书和编纂过程中安排了这段话？它要说服的对象是谁？现实中的军队、仓储、宗族和道路又给这种主张留下了多大余地？
\end{textwindow}

\subsection{简牍、墓葬与编书：思想在传抄和陈列中才成为“文本”}

战国的文本并不等于今日所见的定本。约前300年下葬的郭店楚墓保存儒、道等竹简，\eventanchor{WQ-0300-GUODIAN}{战国·郭店楚简}是理解早期篇章形态的关键物证；楚地贵族与士人以竹简抄写、收藏、传习礼、法和修身知识的网络，见\eventanchor{WQ-0300-CHU-BAMBOO-CIRCULATION}{战国·楚地竹简传习}。它们使楚地而非齐、秦都城的知识生产变得可见，也提醒我们：传世《老子》、儒家篇章或后来的“学派”名称，不能不经比较便当作墓主当日读到的原样全本。一座墓所保存的是被选择、携带和安置的一组文本，不是全国读者的统计样本。

屈原投汨罗江通常系于前278年前后，\eventanchor{WQ-0278-QUYUAN}{战国·屈原投汨罗江的传统系年}；《史记·屈原贾生列传》与《楚辞》使郢都失陷、楚廷迁徙和诗人的政治失意在后世记忆中相互连结。这一传统说明亡国危机如何进入文学与仪式记忆，却不能据以坐实投江的日期、地点和每一段传说细节，更不能将《离骚》直接当作楚廷现场的编年记录。

战国末年，文本还成为政治集团可见的声望。前239年前后，吕不韦主持门客编纂并陈列《吕氏春秋》，见\eventanchor{WQ-0239-LUSHI}{战国·《吕氏春秋》编成}与\eventanchor{WQ-0239-LUSHI-KNOWLEDGE-NETWORK}{战国·吕不韦门客编书}。所谓“悬之国门”的故事，\eventanchor{WQ-0239-LUSHI-BOOK-MARKET}{战国·《吕氏春秋》悬书}，本身已是书籍、门客和相国声望彼此增益的文化记忆；它足以显示咸阳的知识编纂与相权竞争相接，却不能证明一部兼容诸家的书就是秦国官定意识形态，或每一篇都出自吕不韦本人。

同样，李斯前237年的《谏逐客书》，见\eventanchor{WQ-0237-LISI-MEMORIAL-ROADS}{战国·李斯谏逐客}，以道路、货物、兵器和人才说明秦不能轻弃外来资源。它是一篇为具体政治选择服务的政论，能够照亮客卿网络怎样被说成统一战争的条件；不能据其修辞推定所有外来者都受同等欢迎，或把秦廷的每项决定归因于一封奏书。游士能改变可供选择的语言和联盟，不能取消国君、宗室、官僚和战争资源之间的冲突。

\subsection{统一帝国：礼、标准与控制知识的限度}

秦统一后，战国的辩论没有停止，只是被置入更大的行政尺度。文字、车轨和度量衡标准的推行，见\eventref{QIN-0221-STANDARDS}{秦·统一文字与度量衡}；历法与十月岁首的安排，见\eventanchor{QIN-0221-QIN-CALENDAR}{秦·统一时间秩序}。这些措施使赋役、仓储、文书和祭祀较易互相衔接，却不应被写成各地旧习一夕消失。泰山、琅邪等刻石是近同时的帝国自述：\eventanchor{QIN-0219-TAISHAN-INSCRIPTION-LANGUAGE}{秦·泰山刻石}以功德、法令和家族伦理说明统一，\eventanchor{QIN-0219-LANGYA-OCEAN-RITUAL}{秦·琅邪台巡行}则把海祭、刻石和沿海行政连在一起。它们比后出的传记更接近帝国如何希望被看见，却正因如此也带有强烈的宣示目的。

礼仪并未因法令而自动同声。前219年封禅之前，博士围绕礼制形式提出不同方案，见\eventanchor{QIN-0219-TAISHAN-RITUAL-COUNCIL}{秦·泰山礼仪议论}；这提示新皇帝仍须同齐鲁知识传统协商，而非单以武力创造一套无人异议的礼。同期徐福等出海求仙的叙事，见\eventanchor{QIN-0219-XUFU-SEA-EXPEDITION}{秦·徐福求仙}，以及方士承诺失效后逃亡、追责的材料，见\eventanchor{QIN-0212-ALCHEMIST-ACCOUNTABILITY}{秦·方士求仙受追责}，应当同《史记》的成书与政治宗教叙事一起读；它们说明巡行、海洋想象、资源耗费和不确定的技术承诺如何相互缠绕，不能替我们补出可核验的远航图。

前213年的焚书令也不能只用“焚书坑儒”四字结束。令文所述医药、卜筮、种树与博士官藏书的保留，见\eventanchor{QIN-0213-BOOK-BAN-ARCHIVES}{秦·焚书令与官藏}和\eventanchor{QIN-0213-BOOK-BAN-DOCTOR-EXEMPTION}{秦·博士官藏书获保留}，显示秦廷意在压缩借古讽今与私议的空间，同时仍需要行政、技术、礼制和教育知识。它是收拢解释权的尝试，不是知识从社会中完全消失的证明；郡县执行的差异与私下传抄的延续都不能轻率抹去。

前212年前后方士、儒生受处置，后世概括为“坑儒”，见\eventanchor{QIN-0212-SCHOLAR-PUNISHMENT-TRADITION}{秦·方士儒生受处置的叙事}。现场人数、对象和处置细节不足以让我们把汉代以后的定型记忆直接当成秦廷档案。前211年沉璧、陨石和“祖龙死”一类预兆叙事，\eventanchor{QIN-0211-CHENBI-PROPHECY}{秦·沉璧谶言}与\eventanchor{QIN-0211-IMPERIAL-OMENS}{秦·帝国异象}，也更适合说明后人怎样把高压、巡行和帝国末年的不安组织成秦亡故事，而不是当时民意的统计。

约前278年屈原投汨罗江的传统系年，\eventref{WQ-0278-QUYUAN}{战国·屈原投江传统}，同样不能脱离文本的形成史来阅读。《史记·屈原贾生列传》把个人失意、楚国政治和《离骚》式的忠愤连成可感的故事；屈原与楚怀王、顷襄王的基本关联可以进入楚国危机的读法，投江年月、作品归属及传记场景却都有长期争论。它照亮的是后人如何以一位诗人安放失郢前后的楚国记忆，不是为一桩死亡提供无媒介的现场证词。

约前283年齐人王蠋拒降而死的叙事，\eventref{WQ-0283-QI-WANGZHU-DIES}{战国·王蠋拒降}，则把燕军占领齐地时的士人忠义写成一个可传诵的选择。它不能让我们统计齐地有多少人抵抗、谁在何日被迫迁徙；但与田单复齐的叙事并读，能看见亡国与占领如何被后世转化为关于名节、地方服从和复国正当性的语言。思想史在这里不在战场之外，而在幸存者和后来作者为战场赋予何种意义的工作之中。

\begin{sourcenote}[怎样把思想材料放回事件中]
\sourcebook{《孟子》《墨子》《庄子》《荀子》《韩非子》《吕氏春秋》}保存战国论辩和篇章传承，却都有成书、编次和后世注释的问题；它们宜用来分析可供人们调用的政治语言，不宜逐句当作事件现场记录。\sourcebook{《史记》相关世家、列传与《战国策》}能够追踪游士、国君和说辞在后世叙事中的相遇，仍须防止将修辞化的谈话直接复原为决策过程。

郭店楚简、曾侯乙墓器物、中山王器及秦刻石属于不同性质的材料：它们可以使书写、礼乐、工艺和帝国自我表述获得较近时的物证，却各自只覆盖特定墓葬、地点或官方场合。把传世文本、出土材料和各国的战争、财政、行政事件并读，才可看见战国思想既是争论秩序的语言，也是依赖抄写、教育、资助、道路和权力保护才能流动的实践。
\end{sourcenote}
